\subsection{HLS Kernels}
\label{subsec:kernels}
This section is not intended to propose new HLS kernels.
Instead, we explore the efficient ways to reach the maximum available $M$ on FPGAs and implement the standard kernels as a library with parameterized quantization support, which is reusable across different models (e.g., BERT and GPT models implemented in \S\ref{sec:exp}).

%We implement HLS kernels for each operator in the Transformer model.
% https://docs.xilinx.com/r/en-US/ug1483-model-composer-sys-gen-user-guide/Setting-the-HDL-AIE-Block
\noindent\textbf{Linear Operators.}
Linear operators are the key operators in Transformer models because they are ubiquitous and compute-intensive.
There are two types of linear operators in the Transformer layers, activation-weight matrix multiply (A-W GEMM) with bias, and activation-activation (A-A GEMM) matrix multiply.
A-W GEMM includes the projection layers and the linear layers in the FFN, with weights buffered on-chip and activations streamed from an input FIFO;
A-A GEMM includes the two GEMM operators in the SDP attention, with one of the input activations stored in a double buffer and the other one streamed.

We adopt an output-stationary systolic array architecture to implement the A-W and A-A GEMM process engines.
As shown in Figure~\ref{fig:sa}(a), the systolic array is a 2-D array of MAC units of shape $m_1\times m_2$ with FIFOs connecting different MAC units.
The number of MAC units of linear operator $i$ is actually $M_i$ defined in \S\ref{sub:mac_analysis}.
% \yx{really? how about throughput balancing?}
% \yx{the following discussions only apply to A-W GEMM}
We mainly discuss the A-W GEMM in the following, and the idea can also be applied to the A-A GEMM.
The input activation from the previous layer will be first buffered in an activation buffer.
After $m_1$ elements are buffered, they will be streamed into the systolic array together with the $m_2$ weights.
There is also a fully partitioned output buffer on-chip that allows the outputs to directly write back.

\begin{figure}[!htbp]
    \centering
    \includegraphics[width=0.7\linewidth]{figs/systolic_array.pdf}
    \caption{Systolic array and DSP packing. The yellow blocks in the systolic array represent output buffers.}
    % \yx{this figure only applies to A-W GEMM.}
    \label{fig:sa}
\end{figure}

% \yx{the following discussions only apply to A-W GEMM}
Each MAC unit can be implemented with a single DSP block and can provide one-MAC-per-cycle throughput.
Based on the discussion in \S\ref{subsub:memory_packing}, we adopt the W4A8 quantization scheme for our accelerator design, which maximizes the utilization of available resources.
As a result, the matrix multiplications involve either \texttt{int4} by \texttt{int8} or \texttt{int8} by \texttt{int8} operations, which are too large for LUTs, thus relying primarily on DSPs or specialized compute blocks (e.g., AIE~\cite{xilinxAIE}).
% Additionally, in AMD FPGAs, we can pack two integer multipliers into one DSP48E2 if the bit precision is lower than 8, enabling a DSP to provide up to two MACs per cycle, which allows us to use fewer resources to achieve a larger $M$.
% For instance, the DSP48E2 hard blocks on AMD Xilinx UltraScale+ Family FPGAs can support 18-bit by 27-bit multiplication and accumulation~\cite{u280}, allowing to pack two multiplications into a single DSP48E2 slice for a W4A8 quantized model to save DSP resources.
In AMD FPGAs, the DSP48E2 hard blocks can support 18-bit by 27-bit multiplication and accumulation~\cite{u280}, enabling the packing of two multiplications into one slice for a W4A8 quantized model to save DSP resources and achieve a larger $M$.
Figure~\ref{fig:sa}(b) shows the bit packing method for 4-bit by 8-bit integer multiplications.
One activation is filled into the lower 8 bits of the 18-bit DSP input, and two weights are filled into 0-to-3 and 13-to-16 bit positions of the 27-bit DSP input to avoid overlapping results.
Finally, the two multiplication results are extracted by bit-slicing the 45-bit DSP result.
Notice that since the DSP output is wide enough, we can also pack two 8-bit by 8-bit integer multiplications into one DSP slice by further offsetting the second weight and output.
With DSP packing, we can easily double $M$ to achieve higher performance but with much fewer DSPs.

% According to~\cite{xilinxAIE}, one AIE tile is able to perform 128 MAC/cycle and AIE can work under 1 GHz clock frequency.
% To make AIE and DSP comparable, we can scale the AIE clock frequency to 250 MHz. Hence, one AIE tile can provide 512 MAC units. 


\noindent\textbf{Non-Linear Operators.}
Since quantizing non-linear operators can lead to a significant degradation in model accuracy~\cite{sheng2020qbert,xiao2023smoothquant}, and these non-linear operators are not the bottleneck of the design, we directly implement the floating-point version of the operators in HLS.
Specifically, we buffer a row of elements for softmax and LayerNorm functions, which requires conducting reduction along the last dimension.
Consequently, this approach eliminates the need to wait for the complete results for the computation of these non-linear operators and effectively prevents dataflow stalling.

\begin{comment}
\subsubsection{Softmax.}
Softmax function takes a row of numbers and calculates the probability distribution, as depicted in the following equation

\begin{equation}
     \text{softmax}(X)_{ij} = \frac{\mathrm{e}^{x_{ij}}}{\sum_{k=j}^{l} \mathrm{e}^{x_{ik}}}\,. 
\end{equation}

Implementing softmax on FPGA can be expensive due to its computational complexity.
It requires the exponentiation of each input value, followed by a summation of the exponential values.
It needs two passes on the input data to calculate the results.
This exponentiation operation involves floating-point arithmetic, which is computationally expensive on FPGAs.
If the input of softmax is an integer, we can create a look-up table to store all the possible values.
However, since quantizing non-linear operators may significantly drop the model accuracy, and these non-linear operators are not the bottleneck of the design, we choose to leave it as it is.


\subsubsection{Layer Normalization.}
In contrast to batch normalization~\cite{he2016resnet}, layer normalization computes the normalization statistics directly from the sum of the inputs to neurons within a hidden layer~\cite{vaswani2017transformer}, as shown below,

\begin{equation}
    Y = \frac{X-\mathbb{E}[X]}{\sqrt{\mathrm{Var}[X] + \varepsilon}}\gamma + \beta\,.
\end{equation}

$\gamma$ and $\beta$ are fixed after the training process, while the expectation $\mathbb{E}[X]$ and the variance $\mathrm{Var}[X]$ of the input need to be dynamically calculated during inference.
Notice the variance and expectation have the following relation,

\begin{equation}
\mathrm{Var}[X] = \mathbb{E}[X^2] - \mathbb{E}[X]^2\,.
\end{equation}

We can calculate LayerNorm in two passes.
Similarly, if the LayerNorm is quantized, we can reuse the loop-up table for softmax.
Otherwise, we directly calculate the floating point version.

\subsubsection{GELU.}
The Gaussian Error Linear Unit (GELU)~\cite{hendrycks2016gelu} is a non-linear activation function that weights inputs by their percentile, rather than gates inputs by their sign as in ReLUs~\cite{he2016resnet}.
The GELU activation function is $x\Phi(x)$, where $\Phi(x)$ is the standard Gaussian cumulative distribution function.
The full formula is depicted below,

\begin{equation}
\label{equ:GELU}
\mathrm{GELU}\left(x\right) = x\mathbb{P}\left(X\leq{x}\right) = x\Phi\left(x\right) = x \cdot \frac{1}{2}\left[1 + \mathrm{erf}\left(\frac{x}{\sqrt{2}}\right)\right]\,.
\end{equation}

The original GELU function involves the computation of the cumulative distribution function (CDF) of the Gaussian distribution, which can be computationally expensive when implementing the function in hardware.
To address this issue, researchers have proposed an approximate GELU function that can be more easily computed in hardware.
We use a piecewise linear function that closely approximates the original GELU function~\cite{sheng2020qbert}.
\end{comment}